%%%%%%%%%%%%%%%%%%%%%%%%%%%%%%%%%%%%%%%%%
% Uppsala University Assignment Title Page 
% LaTeX Template
% Version 1.0 (27/12/12)
%
% This template has been downloaded from:
% http://www.LaTeXTemplates.com
%
% Original author:
% WikiBooks (http://en.wikibooks.org/wiki/LaTeX/Title_Creation)
% Modified by Elsa Slattegard to fit Uppsala university
% License:
% CC BY-NC-SA 3.0 (http://creativecommons.org/licenses/by-nc-sa/3.0/)

%\title{Title page with logo}
%----------------------------------------------------------------------------------------
%	PACKAGES AND OTHER DOCUMENT CONFIGURATIONS
%----------------------------------------------------------------------------------------

\documentclass[12pt]{report}
\usepackage[T1]{fontenc}
\usepackage[spanish,es-tabla]{babel}
\usepackage[utf8x]{inputenc}
\usepackage{tocloft}
\renewcommand{\cftfigpresnum}{\bfseries Figura~}
\renewcommand{\cftfigaftersnum}{.\hspace{0.5em}}
\renewcommand{\cfttabpresnum}{\bfseries Tabla~}
\renewcommand{\cfttabaftersnum}{.\hspace{0.5em}}
\newlength{\cftmylength}
\settowidth{\cftmylength}{\cftfigpresnum\cftfigaftersnum}
\addtolength{\cftfignumwidth}{\cftmylength}
\setlength{\cftfignumwidth}{2.6cm}
\setlength{\cfttabnumwidth}{2.6cm}
\usepackage{csquotes}
\usepackage{amsmath}
\usepackage{graphicx}
\usepackage{float}
\usepackage{geometry}
\geometry{margin=1in}
\usepackage{hanging}
\usepackage{parskip}
\usepackage{amssymb}
\usepackage{mathabx}
\usepackage{tabularx}
\usepackage{booktabs}
\usepackage{makecell}
\usepackage[backend=biber, style=apa, url=true]{biblatex}
\addbibresource{bibliography.bib}
\usepackage{caption}
\usepackage{multirow}
\captionsetup{labelsep=period}
\usepackage{longtable}
\usepackage{enumitem}
\usepackage[table]{xcolor}
\usepackage{array}
\usepackage{hyperref}
\hypersetup{
    colorlinks,
    linkcolor={black!50!black},
    citecolor={blue!50!black},
    urlcolor={blue!80!black}
}

\makeatletter
\renewenvironment{abstract}{%
    \if@twocolumn%
      \section*{\abstractname}%
    \else %% <- here I've removed \small
      \begin{center}%
        {\bfseries \Large\abstractname\vspace{\z@}}%  %% <- here I've added \Large
      \end{center}%
      \quotation%
    \fi}
    {\if@twocolumn\else\endquotation\fi}
\makeatother


\begin{document}

\begin{titlepage}

\pagenumbering{Alph}
\newcommand{\HRule}{\rule{\linewidth}{0.5mm}} % Defines a new command for the horizontal lines, change thickness here

\center% Center everything on the page
 
%----------------------------------------------------------------------------------------
%	HEADING SECTIONS
%----------------------------------------------------------------------------------------

\textsc{\LARGE Universidad del Valle}\\[0.5cm] % Name of your university/college
\includegraphics[width=3cm]{Recursos/420px-Univalle.svg.png}\\[0.5cm] % Include a department/university logo - this will require the graphicx package
\textsc{\Large Trabajo Profesional}\\[0cm] % Major heading such as course name
\textsc{\large }\\[0cm] % Minor heading such as course title

%----------------------------------------------------------------------------------------
%	TITLE SECTION
%----------------------------------------------------------------------------------------

\HRule\\[0.4cm]
{ \Large \bfseries Implementación y Evaluación Experimental de un Modelo de Lenguaje Basado en Arquitecturas MatMul-Free y Cuantización Ternaria}\\[0cm] % Title of your document
\HRule\\[1.5cm]
 
%----------------------------------------------------------------------------------------
%	AUTHOR SECTION
%----------------------------------------------------------------------------------------

\begin{minipage}{0.6\textwidth}
\begin{center} \large
Juan Camilo Valencia Rivas\\
juan.valencia.rivas@correounivalle.edu.co
\end{center}

\end{minipage}\\[1cm]

\begin{minipage}{0.6\textwidth}
\begin{center} \large
\emph{Director:}\\
Oscar Fernando Bedoya Leiva, PhD.
\end{center}

\end{minipage}\\[1cm]

\begin{minipage}{0.6\textwidth}
\begin{center} \large
\emph{Co-director:}\\
Raúl Ernesto Gutiérrez de Piñerez, PhD.
\end{center}

\end{minipage}\\[1cm]

\vfill

\begin{minipage}{0.8\textwidth}
\begin{center} \large
Universidad del Valle\\
Facultad de Ingeniería\\
Escuela de Ingeniería de Sistemas y Computación\\
Programa Académico de Ingeniería de Sistemas\\
Cali, 2026
\end{center}
\end{minipage}\\


\end{titlepage}

\pagenumbering{roman}

\begin{abstract}
Los Grandes Modelos de Lenguaje (LLMs) basados en arquitecturas Transformer convencionales enfrentan limitaciones estructurales significativas que comprometen su escalabilidad y versatilidad. Estas limitaciones se manifiestan principalmente en dos dimensiones: el costo computacional prohibitivo asociado a las operaciones de multiplicación de matrices (MatMul) y la dependencia crítica de procesos de tokenización que introducen sesgos lingüísticos y obstaculizan el procesamiento eficiente de secuencias no convencionales. Este trabajo de grado aborda estas problemáticas mediante el diseño, implementación y evaluación de una arquitectura innovadora que integra sinérgicamente dos paradigmas emergentes: modelos MatMul-Free y mecanismos de Dynamic Chunking. La arquitectura propuesta, denominada HNetBit, elimina completamente las operaciones MatMul mediante la implementación de pesos ternarios (-1, 0, +1) y operaciones elemento a elemento, reduciendo drásticamente los requerimientos computacionales y energéticos. Paralelamente, incorpora un sistema de chunking dinámico multi-nivel que permite el aprendizaje automático de estrategias de segmentación adaptativa, eliminando la necesidad de tokenización predefinida y facilitando el procesamiento directo de secuencias de bytes crudos. La evaluación comprehensiva del modelo se realizó mediante un benchmark compuesto por texto natural en español, utilizando métricas estandarizadas de precisión y eficiencia computacional. Los resultados obtenidos en trabajos seminales sugirieron que esta integración sinérgica permitió alcanzar un rendimiento competitivo con arquitecturas Transformer tradicionales, ofreciendo a su vez ventajas significativas en eficiencia computacional y capacidad de generalización a través de dominios estructuralmente diversos, estableciendo así las bases para LLMs más sostenibles y versátiles.
\end{abstract}

\vspace{0.5cm}
\noindent\textbf{Palabras clave:} Grandes Modelos de Lenguaje, arquitectura MatMul-Free, cuantización ternaria, chunking dinámico, eficiencia computacional.

\newpage

\begin{center}
    \tableofcontents
\end{center}


\newpage

\begin{center}
    \listoffigures
\end{center}

\newpage

\begin{center}
    \listoftables
\end{center}

\newpage

\renewcommand{\thepage}{\arabic{page}}
\setcounter{page}{1}
\input{Contenido/1. Introducción}

\input{Contenido/2. Marco Referencial}

\input{Contenido/3. Estado del Arte}

\input{Contenido/4. Arquitectura}

\input{Contenido/5. Optimizaciones}

\input{Contenido/6. Modelos y Método Experimental}

\input{Contenido/7. Resultados}

\input{Contenido/8. Discusión}

\input{Contenido/9. Conclusiones}

\input{Contenido/10. Trabajo Futuro}

\clearpage
\newpage
\nocite{*}
\chapter*{Referencias}
\addcontentsline{toc}{section}{\protect\numberline{}Referencias}
\printbibliography[heading=none]

\end{document}